\section{The equations of stellar evolution}

Start of with the mass element

\begin{equation*}
    dm = \rho dV = \rho 4 \pi r^2 dr
\end{equation*}

From the first law of thermodynamics (assuming the conservation of mass with $dm \equiv {\rm const}$) 

\begin{equation*}
    \delta (u dm) = dm \delta u = \delta Q + \delta W
\end{equation*}

where $u$ is the internal energy per unit mass. We can calculate the work $\delta W$

\begin{equation*}
    \delta W = - P delta dV = - P \delta \left(\frac{dV}{dm} dm \right) = - P \delta \left( \frac{1}{\rho} \right) dm
\end{equation*}

Next, note that heat comes from two main sources: release of nuclear energy and the balance of the heat fluxes, going in and out. The rate of nuclear energy per unit mass is $q$ and The heat flowing though the surface is $F(m)$ (units of power), and clearly $F(M) = L$. Thus for heat

\begin{equation*}
    \delta Q = q dm \delta t + F(m) \delta t - F(m + dm) \delta t
\end{equation*}

Using $F(m + dm) = F(m) + \frac{\partial F}{\partial m} dm$, leaves us with

\begin{equation*}
    \partial Q = \left( q - \frac{\partial F}{\partial m} \right) dm \delta t
\end{equation*}

Inserting everything back into firt law we obtain

\begin{equation*}
    dm \delta u + P \delta \left( \frac{1}{\rho} \right) dm = \left( q - \frac{\partial F}{\partial m} \right) dm \delta t
\end{equation*}

It seems that $dm$ cancels and taking the limit $\delta t \rightarrow 0$ we get

\begin{equation*}
    \dot{u} + P \dot{ \left( \frac{1}{\rho} \right)} = q - \frac{\partial F}{\partial m}
\end{equation*}

where the dot is a \textbf{partial} derivative wrt time. In thermal equilibrium there is no change in these quantities wrt time, hence the time derivatives vanish

\begin{equation*}
    q = \frac{dF}{dm} \quad \Rightarrow \quad \int_0^M q dm = \int^M_0 dF = L
\end{equation*}

This integration gives the total power supplied by nuclear processes, denoted by $L_{\rm nuc}$ as \textbf{nuclear luminosity}. Thermal equilibrium implies that the energy created is radiated away at the same rate (No energy build up).

Next we can derive the equations of motion. Start with some mass

\begin{equation*}
    \Delta m = \rho dr dS
\end{equation*}

Any piece of mass feels two opposing forces: gravitational force and force from pressure. Due to spherical symetry we can only consider the radial part of the acceleration

\begin{equation*}
    \ddot{r} \Delta m = - \frac{G m \delta m}{r^2} + P(r) dS - P(r + dr) dS
\end{equation*}

Again, $P(r + dr) = P(r) + \frac{\partial P}{\partial r} dr$, so

\begin{equation*}
    \ddot{r} \Delta m = -\frac{G m \Delta m}{r^2} - \frac{\partial P}{\partial r} \frac{\Delta m}{\rho} \quad \Rightarrow \quad \ddot{r} = \frac{G m}{r^2} - \frac{1}{\rho} \frac{\partial P}{\partial r}
\end{equation*}

Then using $dr = dm/(4 \pi r^2 \rho)$

\begin{equation*}
    \ddot{r} = -\frac{G m}{r^2} - 4 \pi r^2 \frac{\partial P}{\partial m}
\end{equation*}

When acceleration is negligable we get \textbf{hydrostatic equilibrium}

\begin{equation*}
    \frac{dP}{dr} = -\rho \frac{G m}{r^2} \quad \text{or} \quad \frac{dP}{dm} = - \frac{Gm}{4 \pi r^4}
\end{equation*}

We can estimate the pressure by integrating from the center to the outside

\begin{equation*}
    P(M) - P(0) = - \int_0^M \frac{G m dm}{4 \pi r^4} \quad \Rightarrow \quad P_c = \int_0^M \frac{G m dm}{4 \pi r^4} > \int_0^M \frac{G m dm}{4 \pi R^4}
\end{equation*}


\subsection{Virial Theorem}
Multiplying the hydrostatic equilibrium equation by volume $V$

\begin{equation*}
    \int_0^{P(R)} V dP = - \frac{1}{3} \int_0^M \frac{G m dm}{r}
\end{equation*}

But the RHS is just the gravitational potential energy $\Omega$ and we can integrate LHS by parts where boundary terms vanish since $V(0) = 0$ and $P(R) = 0$. Yielding

\begin{equation*}
    -3 \int_0^{V(R)} P dV = \Omega \quad \Leftrightarrow \quad -3 \int^M_0 \frac{P}{\rho} dm = \Omega
\end{equation*}

We can use this result, for ideal gas.

\begin{equation*}
    u = \frac{3}{2} \frac{k T}{m_g} = \frac{3}{2} \frac{P}{\rho}
\end{equation*}

Then combining these equations and noting that $U = \int u dm$ we get

\begin{equation*}
    U = -\frac{1}{2} \Omega
\end{equation*}

We can use these formulas to derive average temperature. For $\Omega = - \alpha \frac{GM^2}{R}$ and $U = \frac{3}{2} \frac{k}{m_g} \bar{T} M$ we obtain

\begin{equation*}
    \bar{T} = \frac{\alpha}{3} \frac{m_g G}{k} \frac{M}{R}
\end{equation*}


\subsection{Total Energy of the star}

